\chapter{En tømrer}

Da trompettonene lød gjennom landet,
skyndte mennene seg av sted. Fra byer, landsbyer og
enslige gårder kom de vandrende mot Jerusalem,
hver med sin stav i hånden, uten å vite hvorfor
Herren hadde kalt dem sammen.

Josef var en gammel tømrer fra Nasaret. Akkurat denne dagen var
han i full sving i verkstedet sitt for å lage bjelker til
et nytt hus som skulle bygges i utkanten av
Nasaret. Nye bjelker måtte hugges til.
Han satte foten mot emnet og løftet øksen.
Slagene kom med jevn rytme --- faste, men aldri hastige.
Flisene spratt til side og samlet seg som lyse spon rundt føttene hans.
Når den grove formen var hogd frem, la han øksen fra seg og
tok frem høvelen. Lange, tynne trespon krøllet
seg sammen og falt lydløst til bakken,
mens den ru overflaten langsomt ble glatt
under de sterke hendene hans.
Luften fyltes av den søtlige duften av ferskt sedertre og eik.


Da tempelets sendebud kom fram til  Josefs verksted, 
la han straks høvelen fra seg der han stod og arbeidet.
Kallet fra tempelet kunne ikke ignoreres. Som de andre enkemennene
tok han vandrestaven sin og dro til møtestedet som sendebudet hadde kunngjort.
Snart var de alle samlet, unge og gamle, hver med en
vandringsstav i hånden og med spørsmål om hvorfor Herren hadde kalt dem.

De fikk høre at de møtte møte opp i tempelet i Jerusalem som
lå 100 romerske mil unna.  Etter fire dagers vandring ankom mennene Jerusalem-
Sammen gikk de opp til tempelet, hvor øverstepresten tok imot dem.
Én etter én samlet han inn stavene deres.
Deretter bar han dem inn i Herrens tempel.
Mens mennene ventet ute i stillhet, gikk han frem for Gud og
ba om at Herren selv måtte vise hvem han hadde utvalgt.

Da bønnen var slutt, kom øverstepresten ut igjen med stavene.
Han ga dem tilbake, én etter én. Alle fulgte spent med,
men ingenting skjedde. Ingen stav skilte seg ut, og intet tegn viste seg.
Til slutt stod bare Josefs stav igjen.

Da øverstepresten rakte den til ham, skjedde det plutselig noe.
En hvit due fløy frem fra staven, steg opp i luften og
slo seg rolig ned på Josefs hode. Et gisp gikk gjennom mengden.
Ingen tvilte lenger på at Herren hadde talt.

Øverstepresten vendte seg mot Josef.
«Herren har valgt deg», sa han.
«Du skal ta Herrens jomfru til deg og ha omsorg for henne.»

Men Josef rygget et skritt tilbake og ristet på hodet.
«Nei», svarte han. «Jeg er en gammel mann, og jeg har allerede sønner.
Hun er så alt for ung. Hvordan skal dette kunne gå?
Folket i Israel vil le av meg.»

Øversteprestens ansikt ble alvorlig.
«Josef», sa han, «frykt Herren din Gud. Husk hva som hendte med
Datan, Abiram og Korah da de gjorde opprør mot Herren.
Jorden åpnet seg og slukte dem. La ikke den samme ulykken ramme
ditt hus ved at du setter deg opp mot Guds vilje.»

Ordene traff Josef. Han kjente fortellingene fra
Skriften og visste hva ulydighet mot Herren kunne føre til.
Frykten fylte ham, men enda sterkere var
vissheten om at han ikke kunne avvise det Gud selv hadde bestemt.

Derfor bøyde han seg for avgjørelsen og tok Maria under sitt vern.
Da han førte henne bort fra tempelet, talte han vennlig til henne:
«Maria, jeg har tatt deg med fra Herrens tempel, og nå fører jeg
deg hjem til mitt hus.»


Sammen med Josef la hun ut på den lange reisen nordover mot Nasaret,
landsbyen der han hadde sitt hjem og sitt verksted.
Veien førte dem ned fra Jerusalems høydedrag,
gjennom det golde landskapet som skrånet mot Jordan-dalen.
Om dagen var luften tørr og varm, mens de kjølige kveldene
fikk dem til å søke ly ved veikroer eller i skyggen av
oliventrær og steingjerder. De møtte handelskaravaner
på vei mot Jerusalem, hyrder som drev saueflokkene
sine mellom beitemarkene og pilegrimer som sang
salmer mens de vandret mot tempelet.

Josef gikk med den lange staven i hånden og passet på at
Maria ikke kom til skade på de steinete stiene.
Selv sa han lite. Han var en mann av få ord
som lot arbeidet og handlingene tale sterkere enn ordene.
Maria fulgte rolig ved siden av ham. 

Etter flere dagers vandring begynte landskapet å forandre seg.
De bratte høydene omkring Jerusalem ble avløst av
Galileas grønnere åser. Vinranker, fikentrær og
olivenlunder bredte seg utover dalsidene, og
små landsbyer lå spredt mellom jordene.
Til slutt fikk de øye på Nasaret, den beskjedne
fjellandsbyen som klamret seg til skråningen.

Der førte Josef henne inn i sitt enkle hus.
Verkstedet lå vegg i vegg med boligen, og duften av nyhogd
sedertre og eik fylte luften. Det var her Maria nå skulle bo.
Slik kom Maria til Josefs hus, ikke som en vanlig brud,
men som den unge jomfruen Gud selv hadde betrodd i den gamle
håndverkerens varetekt.


«Maria, jeg må  dra bort hver dag,
for arbeidet mitt kaller meg. Jeg lever av å bygge hus,
men hver dag når arbeidet er ferdig,
skal jeg vende tilbake til deg», sa Josef.
Han så på henne et øyeblikk før han fortsatte med en roligere stemme.
«Inntil da er du ikke alene. Herren selv skal våke over deg og beskytte deg.»

Noen måneder senere samlet prestene seg igjen til råd i tempelet.
Det gamle forhenget, som skilte den hellige delen av helligdommen fra
Det aller helligste, var blitt slitt av tidens tann.
Stoffet som en gang hadde vært strålende, bar nå preg av mange års
tjeneste. Derfor ble det besluttet at et nytt forheng skulle veves ---
et forheng verdig Herrens hus.

«La oss lage et nytt forheng til Herrens tempel», sa prestene.
«Det skal være vakkert nok til å tjene foran Den Høyestes ansikt.»
Øverstepresten reiste seg.
«Send bud etter de unge jomfruene av Davids stamme», befalte han.
«De som er rene og uklanderlige, skal utføre dette hellige arbeidet.»

Tempelvaktene og budbringerne drog ut.
De lette gjennom Jerusalems gater og landsbyene omkring, og
til slutt fant de sju unge kvinner som oppfylte kravene.
Da de kom tilbake, mintes øverstepresten den unge Maria ---
hun som Herren selv hadde utvalgt og som nå levde under Josefs beskyttelse.
Også hun var av Davids ætt og kjent for sitt rene liv.
«Hent også Maria fra Josef i Nasaret»,  sa han.

Budbringere gikk straks av sted, og allerede vel en uke senere
kom Josef med Maria tilbake.
Hun var stille av natur og vant til å ferdes i de
hellige forgårdene fra årene hun hadde vokst opp der.
Da hun trådte inn gjennom portene, var det som om
hun vendte tilbake til et hjem hun aldri helt hadde forlatt.

De åtte unge kvinnene ble ført frem for prestene.
Foran dem lå nøstene med kostbare tråder som skulle bli til det nye
forhenget –--- glitrende gulltråd, blendende hvit ull, fint lin,
myk silke og garn farget i fiolett, skarlagen og den dypeste purpur,
den kosteligste av alle fargene. Hver tråd var valgt med omhu,
for alt som hørte Herren til, skulle utføres med den største verdighet.

Øverstepresten løftet hendene.
«La loddet avgjøre», sa han. «Herren selv skal vise hvem
som skal arbeide med de forskjellige trådene.»

Det ble kastet lodd, og én etter én fikk jomfruene sin oppgave.
Noen skulle spinne gullet, andre den hvite ullen,
linet, silken eller den fiolette tråden.

Da det siste loddet ble trukket, falt det på Maria.
Hun fikk den skarlagenrøde og den ekte purpurfargede ullen ---
de mest kostbare og praktfulle fargene, fargene som bar
kongelig verdighet og hellighet i seg.

Maria tok imot nøstene med bøyd hode. Hun sa ikke et ord,
men bar dem varsomt med seg hjem til Josefs hus.
Hun betraktet arbeidet som en tjeneste for Gud,
ikke som en ære hun selv hadde fortjent.

På denne tiden ble presten Sakarias rammet av taushet.
Ingen ord kom over leppene hans, slik Herren hadde bestemt.
Inntil stemmen en dag skulle vende tilbake,
overtok Samuel hans tjeneste i tempelet og le
det prestene i de daglige oppgavene.

Hjemme i stillheten satte Maria seg ved rokken.
Sollyset falt inn gjennom døråpningen og traff
den purpurfargede og skarlagenrøde ullen,
så trådene glødet som flammer. Med rolige, øvede hender
begynte hun å spinne. Hjulet dreide jevnt rundt,
og den fine tråden vokste frem mellom fingrene hennes, sterk og jevn.

Ingen visste at mens hun satt bøyd over det hellige arbeidet
og spant trådene til forhenget som en dag skulle henge i Herrens tempel,
var hun selv i ferd med å bli utvalgt til et langt større mysterium
--- et mysterium som ennå var skjult for alle mennesker.


Dagene gikk stille i Josefs hus i Nasaret. Når Josef var borte på arbeid, var huset fylt av en fredelig stillhet. Maria fortsatte trofast oppgaven hun hadde fått av prestene i Jerusalem. Den purpurfargede og skarlagenrøde ullen lå ved rokken hennes, og hver dag spant hun trådene som en dag skulle veves inn i det nye forhenget til Herrens tempel.

En morgen tok hun vannkrukken og gikk ut til landsbyens kilde, slik kvinnene gjorde hver dag. Solen var ennå lav over åsene, og den kjølige morgenluften bar med seg duften av oliventrær og villtimian. Fuglene sang mellom fikentrærne, og vannet rislet klart ned i steinbassenget.

Maria bøyde seg frem for å fylle krukken.

Da lød plutselig en stemme – klar og mild, men uten at hun kunne se hvem som talte.

«Vær hilset, du benådede. Herren er med deg. Velsignet er du blant kvinner.»

Hun stivnet. Hjertet begynte å slå raskere. Forsiktig snudde hun seg mot høyre, så mot venstre. Blikket lette mellom trærne, langs stien og ved kilden. Ingen var der. Bare den stille morgenvinden fikk bladene til å bevege seg.

Stemmen hadde vært virkelig. Det var hun ikke i tvil om.

En dyp ærefrykt grep henne. Hun løftet vannkrukken og skyndte seg hjem med raske skritt. Tankene raste. Hvem hadde talt? Og hvordan kunne noen kjenne henne uten at hun så dem?

Vel hjemme satte hun krukken fra seg ved døren. For å finne ro tok hun frem den purpurfargede ullen fra tempelet. Hun satte seg på stolen ved vinduet, lot rokken begynne å dreie og førte den myke tråden varsomt mellom fingrene. Den jevne rytmen pleide å gi henne fred.

Men denne dagen var stillheten annerledes.

Plutselig ble rommet fylt av et lys som ikke kom fra solen. Foran henne stod en engel, strålende og mektig, men med et ansikt fylt av fred.

Maria ble grepet av frykt. Hun senket blikket, og hendene hennes ble liggende stille over den purpurfargede tråden.

Engelen talte med en stemme som var mild og sterk på samme tid.

«Frykt ikke, Maria. Du har funnet nåde hos Herren, han som hersker over himmel og jord. Etter hans ord skal du bli med barn.»

Ordene fylte henne med undring. Hun forsøkte å fatte det hun nettopp hadde hørt.

Hun løftet blikket forsiktig mot engelen.

«Hvordan kan dette skje?» spurte hun stille.
«Skal jeg bli mor slik andre kvinner blir mødre?»

Engelen smilte. «Nei, Maria. Det skal ikke skje på menneskers vis.
Guds kraft skal komme over deg, og Den Høyestes kraft skal overskygge deg.
Derfor skal barnet som blir født, være hellig.
Han skal kalles Den Høyestes Sønn.
Du skal gi ham navnet Jesus,
for han skal frelse sitt folk fra deres synder.»

Maria satt taus. Hun kjente både ærefrykt og en fred
som langsomt fortrengte frykten. Hun forstod ikke alt
engelen hadde sagt. Hvordan kunne noe slikt skje?
Hvordan kunne Gud velge henne, en ung kvinne fra den lille landsbyen Nasaret?

Likevel kjente hun at Herren, som hadde ledet henne
siden barndommen i tempelet, også ville bære henne gjennom det som nå lå foran.

Hun bøyde hodet.
«Se», sa hun med rolig stemme,
«jeg er Herrens tjenerinne. La det skje med meg slik du har sagt.»

I det samme forsvant lyset, og engelen var borte.

Huset lå igjen i den samme stillheten som før.
Bare rokken stod urørlig ved siden av henne,
og den purpurfargede tråden hvilte mellom fingrene hennes.
Men ingenting var lenger som før. I det stille rommet i
Josefs hus hadde Gud begynt å oppfylle løftene som
profetene gjennom århundrer hadde båret frem.
Den unge kvinnen som satt med trådene til templets forheng,
bar nå selv det levende håpet som en dag skulle åpne veien
mellom Gud og menneskene.
